\section{Induced cubes for a transitive alternating tensor layer}
\label{sec:induced-cube-layer}

\paragraph{Status and provenance.}
D-TD49, reserved question Q2fe, is a new derivation by
\texttt{/root/imprimitive}.  It uses the accepted local character
bounds and the accepted actual sphere-cover genus theorem HD-GC06.
The parent and \texttt{/root/g34\_eightfolds} independently checked
the proposed genus-one interface in discussion.  The complete
author, independent whole-file and root proof readings pass;
their exact scopes and shared contributions are recorded separately.  No support condition
on the external partner, and no twelve-mark condition on a
projected quotient, is imposed.

For a finite-projective-order operator $A$ on a complex vector space
of dimension $d$, write
\[
 \nu(A)=d-\max_{z\in\mathbf C^*}\dim\ker(A-zI).
 \tag{ICL.1}
\]
The normalized support is unchanged by scalar multiplication.
Tensor spectral convolution never decreases it.  A nonidentity
scalar has actual moving rank $d$, although its projective support
is zero.

\subsection*{Actual layer and two precise local inputs}

Let $C\to\mathbf P^1$ be an actual connected finite $J$-Galois
cover, with a finite complex coefficient representation on $V$.
Its genuine generating product-one inertias are $T_1,\ldots,T_r$.
Assume
\[
 e=\dim V,\qquad
       \sum_{i=1}^r\operatorname{rank}(1-T_i)=2e+6.
 \tag{ICL.2}
\]
Suppose there is an actual stable whole tensor layer
\[
 V=\left(\bigotimes_{j=1}^t W_j\right)\otimes R,
 \qquad \dim W_j=a,\qquad
 m\ge10^6,\quad 2\le t\le54m.
 \tag{ICL.3}
\]
Here $R$ is the entire remaining factor, of any positive dimension.
The actual projective image on the selected layer contains the
irreducibly represented normal central product of $t$ quasisimple
components with simple quotients $A_m$.  It permutes their full
matrix algebras transitively.  Thus conjugation gives an actual
homomorphism from $J$ to $S_m\wr S_t$ whose image contains $A_m^t$.
An element in the kernel of the coefficient action fixes every
embedded component matrix algebra pointwise and hence centralizes
every represented local core.  It therefore kills this natural
wreath action, which factors through the actual coefficient image.
Every local cycle product is a compatible normalizer of its
represented irreducible component.  These are hypotheses on the
actual action, not on a proposed independent lift of the components
to $J$.

Fixed component identifications, including odd relabelings when
needed, are chosen once.  For a compatible local operator $A$,
let $y\in S_m$ be the natural permutation supplied by conjugation
on its simple quotient, and let $f(y)$ count its fixed letters.
Assume one of the following two uniform sets of local hypotheses.
\begin{description}
\item[Spin input S.]
For every such operator, after removal of its scalar,
\[
 a\ge m^3,\qquad
 \frac{|\operatorname{tr}(A)|}{a}
       \le\left(\frac{\max\{f(y),1\}}m\right)^3,
 \qquad y\ne1\Longrightarrow
       \frac{|\operatorname{tr}(A)|}{a}\le\frac78.
 \tag{ICL.4}
\]
\item[Strong ordinary input O.]
For every such operator, after removal of its scalar,
\[
 \begin{gathered}
 a\ge m^6,\qquad
 \frac{|\operatorname{tr}(A)|}{a}
 \le\begin{cases}(f(y)/m)^3,&f(y)>0,\\m^{-6},&f(y)=0,
       \end{cases}\\
 y\ne1\Longrightarrow \frac{\nu(A)}a\ge\frac1{m-1}.
 \end{gathered}
 \tag{ICL.5}
\]
\end{description}
All these assertions apply to every compatible cycle product and
every power, with the actual natural permutation of that power.
They are assertions about absolute traces, so scalar phases do not
alter them.  HD-FT23--24 supplies S for every genuine spin
irreducible, including compatible outer operators.  HD-TW02,
HD-FT02 and HD-FT03 supply the character part of O for ordinary
extendible constituents of degree at least $m^{120}$ and for all
ordinary split constituents in this range.  HD-FT20 supplies its
support floor.  The exact previously reviewed source scopes are
listed after the proof.

\begin{theorem}[D-TD49]
Under (ICL.2)--(ICL.3), either local input S or local input O
is impossible.  The external partner $R$, all scalar phases,
all original coefficient fields, and all original inertia orders
are unrestricted.  In particular the statement does not assume
that the selected layer is itself a sixfold supplier.
\end{theorem}

\subsection*{Weighted traces and an induced permutation comparator}

\begin{proof}
Let $g$ be any actual finite-order coefficient operator, and let
$\pi$ be its permutation of the $t$ component algebras.  Write
$c$ for the number of fixed components and
$I=t-\#\{\text{cycles of }\pi\}$.
Tensor-factor conjugation puts a weighted cycle of length $\ell$
into cyclic rotation with one twist $A_C$, its local cycle
product.  The elementary trace contraction around each cycle
then gives
\[
 \frac{|\operatorname{tr}(g)|}{e}
       \le a^{-I}\prod_{C}\frac{|\operatorname{tr}(A_C)|}{a}.
 \tag{ICL.6}
\]
Here and below local representatives are normalized to finite
order, and the whole partner contributes at most one to the
normalized absolute trace.

For precision, this normalization loses no original power or
scalar.  A power killing the component permutation is an ordinary
tensor product of local cycle products and a partner power.
Since a further power of $g$ is the identity, each of these
factors has a scalar power.  They are semisimple and can be
rescaled to finite order.  After the displayed cycle conjugation,
each weighted cycle and the partner can therefore be rescaled
to finite order.  The remaining overall scalar is a root of
unity, because $g$ itself has finite order.  Formula (ICL.6) is
thus an equality on the selected layer times a partner factor
of modulus at most one.  The argument applies separately to
every actual power of $g$; it does not replace the original
order by the natural permutation order.

If $c>0$, the product of the $c$ fixed-component ratios, all in
$[0,1]$, is at most their arithmetic mean.  The nonfixed-cycle
ratios are at most one.  Moreover
\[
              a^{-I}\le\frac ct\qquad(c>0).
 \tag{ICL.7}
\]
Indeed this is equality if $\pi=1$.  Otherwise let $u=t-c\ge2$.
Every nontrivial cycle contributes at least half its length to
$I$, so $I\ge u/2$.  Since $a\ge4$,
$a^I\ge2^u\ge u+1\ge1+u/c=t/c$.
If $c=0$, the same index estimate gives $I\ge t/2$.
Consequently
\[
 \frac{|\operatorname{tr}(g)|}{e}
 \le\begin{cases}
 \displaystyle\frac1t\sum_{j:\pi j=j}
              \frac{|\operatorname{tr}(A_j)|}{a},&c>0,\\
 a^{-t/2},&c=0.
 \end{cases}
 \tag{ICL.8}
\]

Let $\Omega$ be the disjoint union, over the $t$ blocks, of the
ordered cubes of their $m$ natural letters, allowing repetitions.
This is an actual $J$-set of size $D=tm^3$.  Its fixed-point
function is
\[
             F_\Omega(g)=\sum_{j:\pi j=j}f(y_j)^3.
 \tag{ICL.9}
\]
By (ICL.4)--(ICL.5) and (ICL.8), for every original power one has
\[
 \frac{|\operatorname{tr}(g)|}{e}
 \le\frac{F_\Omega(g)}D+\epsilon,
 \qquad
 \epsilon=\begin{cases}m^{-3},&\mathrm S,\\m^{-6},&\mathrm O.
 \end{cases}
 \tag{ICL.10}
\]
For $c>0$, each local error is at most the stated $\epsilon$,
and the coefficient $c/t$ is at most one.  For $c=0$, use
$a^{-t/2}\le a^{-1}\le\epsilon$.
If the natural wreath image of $g$ is the identity, then
$F_\Omega(g)=D$, and (ICL.10) holds with error exactly zero.
This remains so when the whole partner or the actual scalar
part of $g$ is nontrivial.

\subsection*{Counting only actual entries that can bear an error}

Let $k$ count those original genuine inertias whose natural
wreath image is nonidentity.  For S, a component-moving inertia
has normalized trace at most $1/a\le7/8$ by (ICL.6).
If all components are fixed, a nontrivial natural image has at
least one local nonidentity permutation, so (ICL.4) again gives
whole trace at most $7/8$.
For a unitary operator $T$, the triangle inequality applied to
an eigenspace of maximum multiplicity gives
$\nu(T)/e\ge(1-|\operatorname{tr}(T)|/e)/2$.
Each of these $k$ entries therefore costs at least $e/16$.
Since $e\ge a^t\ge m^6>96$, (ICL.2) implies
\[
                         k\le32\qquad(\mathrm S).
 \tag{ICL.11}
\]

For O, a component-moving inertia has normalized projective
support at least $1/4$, by the accepted weighted-cycle bound
used in HD-FT18--21.  One can see its precise scope directly:
in a diagonal twist basis, only the $a$ constant words have
cyclic orbit size one; every other word orbit has size at least
two and contributes at most one to any target eigenspace.
Thus a moved cycle of length $\ell\ge2$ has maximum multiplicity
at most $(a^\ell+a)/2\le3a^\ell/4$.
Tensoring any other cycles and the whole partner preserves the
normalized support floor.  If all components are fixed, a
nonidentity natural image has a local support at least
$1/(m-1)$ by (ICL.5); spectral convolution again preserves it.
Thus every counted entry costs at least $e/(m-1)$.
Using $e\ge m^{12}$ in (ICL.2),
\[
 k\le\left(2+\frac6e\right)(m-1)<2m-1,
 \qquad\text{hence }k\le2m\qquad(\mathrm O).
 \tag{ICL.12}
\]
The much weaker last bound is sufficient.  No uniform $7/8$
character bound is asserted for the ordinary case; near-identity
ordinary permutations are the reason for the separate count.

Let $L_i$ be the actual order of $T_i$.  Fourier inversion at
any target eigenvalue and (ICL.10), averaged over all $L_i$
powers, give
\[
 \frac{\nu(T_i)}e
 \ge\frac{\operatorname{ind}_\Omega(T_i)}D-\epsilon
 \quad\text{for a counted entry},
 \qquad
 \sum_i\frac{\operatorname{rank}(1-T_i)}e
 \ge\frac1D\sum_i\operatorname{ind}_\Omega(T_i)-k\epsilon.
 \tag{ICL.13}
\]
The natural order divides $L_i$, so averaging its fixed-point
function over $L_i$ powers gives exactly its number of cycles.
If the original natural image is identity, all its powers have
zero error and zero index; its actual rank is merely retained
as a nonnegative term.  There is no assertion that twelve
original marks remain nontrivial in this quotient.

\subsection*{Five actual covers and every component-base genus}

The action on $\Omega$ has exactly five orbits: the diagonal,
three specified equality patterns with exactly two distinct
letters, and the ordered triples of distinct letters.  The
inner group $A_m^t$ is three-transitive inside each block and
the actual component action is transitive, proving this orbit
claim without a stable individual-component quotient of $J$.
Let $Y_1,Y_2,Y_3$ denote the corresponding actual connected
quotient curves; the three equality-pattern curves are mutually
isomorphic over the sphere by coordinate relabeling.
Riemann--Hurwitz on their five orbits yields
\[
 \sum_i\operatorname{ind}_\Omega(T_i)
 =2D-10+2\bigl(g(Y_1)+3g(Y_2)+g(Y_3)\bigr).
 \tag{ICL.14}
\]

Let $H\le J$ stabilize one component, let $X=C/H$, and put
$h=g(X)$.  The actual map $H\to S_m$ supplied by conjugation
on that component has image $P=A_m$ or $S_m$, because it
contains $A_m$.  Its kernel $K$ gives an actual connected
$P$-Galois cover $C/K\to X$.  The ordered-distinct-triple
quotient of this cover is exactly $Y_3\to X$, of degree
$(m)_3=m(m-1)(m-2)$.  In particular $X$ is not assumed rational.
We prove the common lower bound
\[
                         g(Y_3)>\frac{m^2}{48}-5.
 \tag{ICL.15}
\]

If $h=0$, the accepted HD-GC06 theorem applies to this actual
$A_m/S_m$ sphere cover: its unordered-three-subset quotient $Z$
has genus $g(Z)>m^2/288$.  The natural forget-order map
$Y_3\to Z$ has degree six.  This is true also for $A_m$,
since a permutation of three selected letters can be accompanied
by a permutation of matching parity on the remaining letters.
Hence $g(Y_3)\ge6g(Z)-5>m^2/48-5$.

If $h\ge2$, Riemann--Hurwitz gives
$g(Y_3)\ge1+(m)_3(h-1)\ge1+(m)_3$, which is stronger.

Finally suppose $h=1$.  The $P$-Galois cover $C/K\to X$ must
ramify: a connected unramified cover of a complex elliptic curve
has abelian monodromy, whereas $P$ contains the nonabelian
simple group $A_m$.  Choose a genuine branch permutation
$y\ne1$.  It fixes at most $m-2$ natural letters, and thus at
most $(m-2)(m-3)(m-4)$ ordered distinct triples.  A permutation
has index at least half its number of moved points.  Therefore
this one branch contributes at least
\[
 \frac{(m)_3-(m-2)(m-3)(m-4)}2=3(m-2)^2
 \tag{ICL.16}
\]
to the ramification of $Y_3\to X$.  It follows that
$g(Y_3)\ge1+3(m-2)^2/2$, again stronger than (ICL.15).
This proves (ICL.15) in every possible component-base genus.

\subsection*{The uniform numerical contradiction}

Since all quotient genera are nonnegative, (ICL.14)--(ICL.15)
give
\[
 \frac1D\sum_i\operatorname{ind}_\Omega(T_i)
 >2+\frac1{24tm}-\frac{20}{tm^3}.
 \tag{ICL.17}
\]
For S, (ICL.11) gives $k\epsilon\le32/m^3$.
For O, (ICL.12) gives $k\epsilon\le2/m^5\le32/m^3$.
Thus (ICL.3), (ICL.13), and (ICL.17) imply in either case
\[
 \begin{split}
 \frac1e\sum_i\operatorname{rank}(1-T_i)
 &>2+\frac1{1296m^2}-\frac{42}{m^3}\\
 &>2+\frac1{2592m^2}>2+\frac6e.
 \end{split}
 \tag{ICL.18}
\]
The first line uses $t\le54m$ and $t\ge2$.
The second line follows from the exact integer comparisons
$m>42\cdot2592=108864$ and
$m^4>6\cdot2592=15552$, together with $e\ge m^6$.
They hold for every $m\ge10^6$.
This contradicts (ICL.2) and proves D-TD49.
\end{proof}

\subsection*{Precisely proved near-maximum applications}

For this application, additionally assume that the coefficient
image is faithful and primitive and that $r\ge12$ in (ICL.2),
as in the accepted primitive rank-nine sixfold source interface.
These additional hypotheses are not needed for the standalone
D-TD49 theorem.  Fix the same complete D-TD12 recursion throughout.
Let $n$ be
the maximum alternating natural degree over all of its layers,
and suppose $n\ge N_*$, where
\[
 M_*=1+241\cdot240!\cdot2^{239},\qquad F_*=240M_*+2,
 \qquad N_*=[100(1+M_*)F_*^{240}]^2.
 \tag{ICL.19}
\]
Assume an actual alternating layer in the complement has
$m\ge\lceil(n+120)/108\rceil$, as supplied in its exact scope
by D-TD44.  The accepted general orbit bound D-TD16, using the
entire same-recursion support floor $1/(n-1)$, gives
\[
 t\le\left\lfloor\frac{n-1}2\right\rfloor<54m,
 \qquad
 N_*\ge10000\cdot241^{482},\quad m>241^{482}>10^6.
 \tag{ICL.20}
\]
The middle inequality follows from $1+M_*\ge241$ and
$F_*\ge241$.  The last follows from the assumed lower bound
on $m$; no new cutoff is being introduced by this observation.

If $t\ge2$, D-TD49 excludes every genuine spin layer by S,
every ordinary split layer by O, and every ordinary extendible
layer of component degree at least $m^{120}$ by O.
For the latter case the accepted effective MSP consequence is
$|\chi|/a\le(f/m)^3$ for $f>0$ and $|\chi|/a<m^{-6}$ for
$f=0$, once $m\ge4000$.  Its finite $S_m$ extension supplies
all compatible cycle products.  In the split case, fixed odd
component relabelings identify the chosen constituent, and
every compatible cycle product is inner; HD-FT03 supplies
both $a\ge m^{120}$ and exactly the required trace bounds.
HD-FT20 supplies the full compatible normalizer support floor
in both ordinary cases.  Perfectness and Schur's lemma ensure
that all residual discrepancies are scalar, as in those
accepted interfaces.

If $t=1$, the same three classes are already excluded, with
arbitrary entire partners, by HD-FT28, HD-FT03, and HD-FT02,
respectively.  Thus no new singleton claim is needed.
The arbitrary-shape degree argument in HD-FT02 says that, at
$m\ge241^{482}$, an ordinary extendible partition whose row
and column tails are both at least $241$ has degree at least
$m^{120}$.  Consequently a surviving layer in this near-maximum
range must be ordinary extendible with minimum row/column tail
at most $240$.  A trivial component is impossible because its
represented quasisimple action is faithful and nontrivial.

This is a strict layer-type reduction.  It does not yet prove
that the surviving small-tail layer is standard or a singleton.
In particular it neither replaces $1/(n-1)$ by $1/(m-1)$ on
the entire external partner nor assigns the sixfold cohomology
budget to that partner.  The original cover and all its powers
were retained throughout.

\paragraph{Known inputs and documentary scope.}
The only genus estimate invoked is the fully reviewed HD-GC06
for actual $A_m/S_m$ sphere covers at $m\ge10^6$; the positive
component-base-genus cases above are new elementary additions.
The published spin inputs and exact cover normalization are
those already read and located in HD-FT23--24:
Guralnick--Tiep Lemma~2.23, pp.~621--622, and the complex spin
minimum recorded in its Proposition~3.4 proof, p.~626, together
with the reviewed Kleshchev--Tiep input.\cite{HD:GT2012,HD:KT2004}
The ordinary effective input is the previously fully reviewed
MSP theorem and normalization used in HD-TW02, not a new proof
of that published result.\cite{HD:MSP2007,HD:MullerSLC2006}
The split formula and its full effective reduction are exactly
HD-FT03.  The support input is HD-FT20, using the exact reviewed
Guralnick--Saxl Lemma~6.1.\cite{BI:GS}
The adjacent source/status record binds every accepted dependency
and distinguishes fresh reading of these written interfaces from
earlier primary-source audits.
